% ============================================================
% PREAMBLE.TEX
% Sistema Adaptativo de Classificação Multicritério
% de Fluidos de Corte
%
% FINALIDADE
%
% Centralizar:
%
% - codificação;
% - idioma;
% - matemática;
% - tabelas;
% - figuras;
% - caminhos gráficos;
% - configurações gerais da apresentação.
%
% Este arquivo faz parte apenas da BASE do projeto.
%
% Não deve conter:
%
% - resultados científicos;
% - ranking;
% - pesos;
% - conclusões;
% - conteúdo específico dos slides.
%
% ============================================================


% ============================================================
% CODIFICAÇÃO E IDIOMA
% ============================================================

\usepackage[T1]{fontenc}
\usepackage[utf8]{inputenc}
\usepackage[brazil]{babel}


% ============================================================
% MATEMÁTICA
% ============================================================

\usepackage{amsmath}
\usepackage{amssymb}
\usepackage{mathtools}


% ============================================================
% FIGURAS
% ============================================================

\usepackage{graphicx}


% ============================================================
% CAMINHOS DE FIGURAS
% ============================================================
%
% Como main.tex está localizado em:
%
% templates/presentation/latex/
%
% os caminhos abaixo serão relativos a esse diretório.
%
% assets/
%     logos e materiais externos confirmados.
%
% figures/
%     figuras utilizadas na apresentação.
%
% ============================================================

\graphicspath{
    {assets/}
    {figures/}
}


% ============================================================
% TABELAS
% ============================================================

\usepackage{booktabs}
\usepackage{array}
\usepackage{tabularx}


% ============================================================
% COLUNAS E ALINHAMENTO
% ============================================================
%
% O ambiente columns utilizado nos slides já pertence ao Beamer.
%
% tabularx é carregado acima apenas para permitir tabelas que
% precisem se ajustar automaticamente à largura disponível.
%
% ============================================================


% ============================================================
% MICROTIPOGRAFIA
% ============================================================
%
% Melhora pequenas características tipográficas quando
% suportado pelo compilador.
%
% Não modifica resultados ou conteúdo científico.
%
% ============================================================

\usepackage{microtype}


% ============================================================
% CONFIGURAÇÕES MATEMÁTICAS
% ============================================================
%
% Evita numeração automática de equações na apresentação.
%
% As equações utilizadas nos slides têm função explicativa e,
% normalmente, não precisam de numeração.
%
% ============================================================

\setbeamertemplate{theorems}[numbered]


% ============================================================
% NAVEGAÇÃO
% ============================================================
%
% Os símbolos padrão de navegação do Beamer geralmente ocupam
% espaço visual sem acrescentar informação à apresentação.
%
% Por isso são ocultados nesta base.
%
% ============================================================

\setbeamertemplate{navigation symbols}{}


% ============================================================
% NOTAS
% ============================================================
%
% A apresentação poderá futuramente utilizar notas do Beamer,
% mas nenhuma configuração específica será ativada nesta fase.
%
% Caso os integrantes decidam utilizar:
%
% \note{...}
%
% deverão definir posteriormente como essas notas serão
% compiladas.
%
% ============================================================


% ============================================================
% REFERÊNCIAS
% ============================================================
%
% Nenhum pacote bibliográfico adicional é carregado agora.
%
% Motivo:
%
% - ainda não sabemos como os integrantes irão gerenciar as
%   referências da apresentação;
%
% - nenhuma referência científica foi selecionada ainda;
%
% - não queremos impor BibTeX/Biber sem necessidade.
%
% Se a apresentação final utilizar referências bibliográficas,
% a implementação correspondente será feita pelos integrantes.
%
% ============================================================


% ============================================================
% CORES
% ============================================================
%
% Não definir cores neste arquivo.
%
% As cores visuais da apresentação serão centralizadas em:
%
% theme.tex
%
% Isso evita duplicação de configurações.
%
% ============================================================


% ============================================================
% FONTES
% ============================================================
%
% A apresentação utilizará inicialmente as fontes padrão do
% Beamer.
%
% Não estamos impondo uma fonte institucional nesta fase.
%
% Caso a professora forneça um template visual ou orientação
% tipográfica, ela deverá substituir esta configuração genérica.
%
% ============================================================


% ============================================================
% HIPERLINKS
% ============================================================
%
% O Beamer já fornece suporte interno ao hyperref.
%
% Portanto, não é necessário carregar \usepackage{hyperref}
% novamente.
%
% Configurações específicas de links poderão ser adicionadas
% futuramente se forem necessárias.
%
% ============================================================


% ============================================================
% LEGENDAS
% ============================================================
%
% O Beamer já possui mecanismos próprios para figuras e tabelas.
%
% Evitamos carregar o pacote caption aqui porque ele pode gerar
% conflitos de formatação com o Beamer.
%
% ============================================================


% ============================================================
% FLOATS
% ============================================================
%
% Em apresentações Beamer, figuras e tabelas normalmente não
% precisam de mecanismos tradicionais de posicionamento como:
%
% [htbp]
% [H]
%
% O conteúdo é organizado diretamente dentro dos frames.
%
% Portanto, o pacote float não é necessário nesta base.
%
% ============================================================


% ============================================================
% CAMINHOS ABSOLUTOS
% ============================================================
%
% Nunca utilizar:
%
% C:/Users/Nome/Desktop/...
%
% nos arquivos LaTeX.
%
% Todo recurso deve utilizar caminho relativo ao repositório.
%
% ============================================================


% ============================================================
% ASSETS
% ============================================================
%
% Materiais externos confirmados poderão ser colocados em:
%
% assets/
%
% Exemplos futuros:
%
% assets/unesp.jpg
% assets/estat.png
%
% Não é necessário criar ou adicionar imagens agora.
%
% ============================================================


% ============================================================
% FIGURAS CIENTÍFICAS
% ============================================================
%
% As figuras geradas posteriormente pelos integrantes poderão
% ser disponibilizadas em:
%
% figures/
%
% Exemplos futuros:
%
% figures/pesos_critic.pdf
% figures/ranking_topsis.pdf
% figures/pareto.pdf
% figures/robustez.pdf
%
% Esses nomes são apenas exemplos conceituais.
%
% Não criar esses arquivos agora.
%
% ============================================================


% ============================================================
% FORMATOS RECOMENDADOS
% ============================================================
%
% Para gráficos científicos:
%
% preferencialmente:
%
% PDF
%
% quando forem gráficos vetoriais.
%
% Para imagens externas:
%
% PNG
% JPG
% PDF
%
% conforme o material original.
%
% ============================================================


% ============================================================
% TAMANHO DAS FIGURAS
% ============================================================
%
% Preferir dimensões relativas:
%
% \includegraphics[width=0.90\textwidth]{...}
%
% ou:
%
% \includegraphics[width=\columnwidth]{...}
%
% evitando tamanhos absolutos desnecessários.
%
% ============================================================


% ============================================================
% EQUAÇÕES
% ============================================================
%
% Fórmulas metodológicas centrais poderão aparecer nos slides,
% por exemplo:
%
% CRITIC:
%
% C_j =
% \sigma_j
% \sum_k (1-r_{jk})
%
% TOPSIS:
%
% C_i =
% S_i^- / (S_i^+ + S_i^-)
%
% Este arquivo apenas fornece os pacotes necessários para que
% essas equações possam ser compiladas.
%
% ============================================================


% ============================================================
% TABELAS GRANDES
% ============================================================
%
% Evitar tabelas extensas em slides.
%
% Caso uma tabela não caiba em um frame, os integrantes deverão
% avaliar:
%
% - reduzir o número de colunas;
% - mostrar somente as linhas importantes;
% - transformar a informação em gráfico;
% - utilizar slide adicional.
%
% Não reduzir excessivamente a fonte apenas para fazer uma
% tabela caber.
%
% ============================================================


% ============================================================
% RESPONSABILIDADE DOS INTEGRANTES
% ============================================================
%
% Durante a execução, os integrantes poderão acrescentar pacotes
% necessários para necessidades reais.
%
% Exemplos:
%
% - pacote para diagramas;
% - pacote bibliográfico;
% - pacote para tabelas específicas;
% - suporte adicional a código.
%
% Entretanto, cada dependência nova deve possuir finalidade
% clara.
%
% Evitar adicionar pacotes apenas por precaução.
%
% ============================================================


% ============================================================
% MATERIAL DA PROFESSORA
% ============================================================
%
% Se a professora fornecer posteriormente um template Beamer,
% um arquivo .sty, uma identidade visual ou regras específicas
% de apresentação:
%
% material oficial
% >
% configuração genérica desta base
%
% As alterações deverão ser realizadas mantendo uma cópia ou
% referência do material original.
%
% ============================================================


% ============================================================
% COMPATIBILIDADE
% ============================================================
%
% A configuração foi mantida simples para facilitar a
% compilação por diferentes ambientes LaTeX.
%
% A validação final da compilação será responsabilidade dos
% integrantes durante a execução do projeto.
%
% ============================================================


% ============================================================
% ESTADO ATUAL
% ============================================================
%
% PREAMBLE_STATUS = BASE_PREPARED
%
% SCIENTIFIC_RESULTS = NOT_EXECUTED
% REFERENCES_CONFIGURED = FALSE
% OFFICIAL_THEME_AVAILABLE = FALSE
%
% ============================================================